% ==========================================================
% DeepSeis: Makine Öğrenmesi Tabanlı Sismik Anomali Tespiti
% ve Olay Çevresi Davranış Analizi
% BİTİRME TEZİ
% ==========================================================
\documentclass[12pt,a4paper,oneside]{book}

% Türkçe dil desteği
\usepackage[utf8]{inputenc}
\usepackage[T1]{fontenc}
\usepackage[turkish]{babel}

% Sayfa düzeni
\usepackage[margin=2.8cm]{geometry}

% Matematik ve semboller
\usepackage{amsmath}
\usepackage{amsfonts}
\usepackage{amssymb}

% Görseller ve tablolar
\usepackage{graphicx}
\usepackage[export]{adjustbox}
\usepackage{float}
\usepackage{booktabs}
\usepackage{array}
\usepackage{longtable}
\usepackage{multirow}
\usepackage{tabularx}
\usepackage{caption}

% Referanslar ve linkler
\usepackage{hyperref}
\usepackage[numbers,sort&compress]{natbib}

% Başlık ve sayfa numaralandırma
\usepackage{fancyhdr}
\usepackage{titlesec}

% Satır aralığı
\usepackage{setspace}
\onehalfspacing
\raggedbottom
\emergencystretch=2em
\setlength{\parindent}{1.25cm}
\setlength{\parskip}{0.15em}

% Şekil/tablo aralıkları
\captionsetup{font=small,labelfont=bf,skip=6pt}
\setlength{\textfloatsep}{12pt plus 2pt minus 2pt}
\setlength{\floatsep}{10pt plus 2pt minus 2pt}
\setlength{\intextsep}{10pt plus 2pt minus 2pt}
\renewcommand{\floatpagefraction}{0.75}
\renewcommand{\textfraction}{0.12}
\renewcommand{\topfraction}{0.88}
\renewcommand{\bottomfraction}{0.80}

% Gorsel ve tablo tasma kontrolu
\newcommand{\fitfig}[2][]{%
    \makebox[\textwidth][c]{%
        \includegraphics[max width=\textwidth,max height=0.64\textheight,keepaspectratio,#1]{#2}%
    }%
}

% Listeler
\usepackage{enumitem}
\setlist{itemsep=2pt, topsep=4pt, parsep=0pt, partopsep=0pt}

% Kod blokları
\usepackage{listings}
\usepackage{xcolor}

% Algoritma
\usepackage{algorithm}
\usepackage{algorithmic}

% PDF metadata
\hypersetup{
    colorlinks=true,
    linkcolor=black,
    filecolor=magenta,
    urlcolor=blue,
    citecolor=black,
    pdftitle={DeepSeis - Sismik Anomali Tespiti},
    pdfauthor={Can Ahmedi Yaşar PARLAK}
}

% Sayfa başlığı ayarları - sayfa numarası alt ortada
\pagestyle{fancy}
\fancyhf{}
\fancyfoot[C]{\thepage}
\renewcommand{\headrulewidth}{0pt}
\renewcommand{\footrulewidth}{0pt}

% Bölüm başlangıç sayfaları için de aynı stil
\fancypagestyle{plain}{
    \fancyhf{}
    \fancyfoot[C]{\thepage}
    \renewcommand{\headrulewidth}{0pt}
    \renewcommand{\footrulewidth}{0pt}
}

% Bölüm başlık formatı
\titleformat{\chapter}[display]
{\normalfont\Large\bfseries}{\chaptertitlename\ \thechapter}{20pt}{\LARGE}
\titlespacing*{\chapter}{0pt}{-20pt}{24pt}

% ==========================================================
\begin{document}

% -------------------- KAPAK SAYFASI --------------------
\begin{titlepage}
    \centering
    \vspace*{0.3cm}
    
    % Üniversite Logosu
    \fitfig[max width=3.5cm,max height=4.8cm]{figures/RTEU-EN LOGO.png}
    
    \vspace{0.6cm}
    
    {\Large\bfseries RECEP TAYYİP ERDOĞAN ÜNİVERSİTESİ\par}
    {\large MÜHENDİSLİK VE MİMARLIK FAKÜLTESİ\par}
    {\large BİLGİSAYAR MÜHENDİSLİĞİ BÖLÜMÜ\par}
    
    \vspace{1.2cm}
    
    {\LARGE\bfseries Makine Öğrenmesi Tabanlı Sismik\\
    Anomali Tespiti ve Olay Çevresi\\
    Davranış Analizi\par}
    
    \vspace{0.4cm}
    {\Large\bfseries ``DeepSeis''\par}
    
    \vspace{1.2cm}
    {\Large\bfseries BİTİRME TEZİ\par}
    
    \vspace{1.2cm}
    {\Large Can Ahmedi Yaşar PARLAK\par}
    
    \vfill
    
    {\large OCAK 2026\par}
    {\large RİZE\par}
\end{titlepage}

% -------------------- İÇ KAPAK SAYFASI --------------------
\newpage
\thispagestyle{empty}
\begin{center}
    \vspace*{1cm}
    
    {\Large\bfseries Makine Öğrenmesi Tabanlı Sismik\\
    Anomali Tespiti ve Olay Çevresi\\
    Davranış Analizi\par}
    
    \vspace{0.5cm}
    {\Large\bfseries ``DeepSeis''\par}
    
    \vspace{2cm}
    {\Large Can Ahmedi Yaşar PARLAK\par}
    
    \vspace{3cm}
    
    \begin{adjustbox}{center,max width=\textwidth}
    \begin{tabular}{ll}
        \textbf{Jüri Üyeleri} & \\
        Danışman: & Dr. Öğr. Üyesi Tuğçem PARTAL \\
        Üye: & \\
        Üye: & \\
    \end{tabular}
    \end{adjustbox}
    
    \vspace{2cm}
    
    \textbf{Bölüm Başkanı:} Dr. Öğr. Üyesi Yasin YALÇIN
    
    \vfill
    
    {\large OCAK 2026\par}
    {\large RİZE\par}
\end{center}

% -------------------- ÖNSÖZ --------------------
\newpage
\pagenumbering{Roman}
\setcounter{page}{3}
\chapter*{ÖNSÖZ}
\addcontentsline{toc}{chapter}{ÖNSÖZ}

Depremler, yüksek yıkıcılık potansiyeli nedeniyle hem insan hayatı hem de ekonomik açıdan kritik öneme sahip doğal afetlerdir. Türkiye, Alp-Himalaya deprem kuşağı üzerinde yer alması nedeniyle tarih boyunca sayısız yıkıcı depreme maruz kalmıştır. Bu gerçeklik, deprem bilimi alanındaki araştırmaları ülkemiz için hayati bir öncelik haline getirmektedir.

Bu tez çalışması, yapay zekâ ve makine öğrenmesi teknolojilerinin sismik veri analizine entegrasyonunu araştırmaktadır. Çalışmanın temel motivasyonu, STFT tabanlı zaman-frekans öznitelikleri ve sınıflandırma modelleri ile sismik sinyallerdeki normal/anomali ayrımını değerlendirmek ve bu ayrımın deprem olay çevresiyle ilişkisini nicel olarak analiz etmektir.

Tezin hazırlanması sürecinde değerli bilgi birikimi, rehberliği ve yönlendirmeleri ile akademik gelişimime büyük katkı sağlayan danışman hocam Dr. Öğr. Üyesi Tuğçem Partal'a en içten teşekkürlerimi sunarım.

Bu çalışma, deterministik bir deprem tahmini iddiası taşımamakta; bunun yerine sismik sinyallerdeki istatistiksel değişimlerin makine öğrenmesi yöntemleriyle ölçülebilir ve tekrarlanabilir biçimde analiz edilebileceğini ortaya koymayı hedeflemektedir.

\vspace{2cm}
\begin{flushright}
    Can Ahmedi Yaşar PARLAK\\
    Ocak 2026, Rize
\end{flushright}

% -------------------- ÖZET --------------------
\chapter*{ÖZET}
\addcontentsline{toc}{chapter}{ÖZET}

Bu tez çalışması, sürekli sismik kayıtlar üzerinde makine öğrenmesi tabanlı anomali analizi gerçekleştirmeyi hedeflemektedir. Çalışmada MiniSEED formatındaki sismik kayıtlar kısa zaman pencerelerine ayrılmış, her pencere kısa süreli Fourier dönüşümü (STFT) ile zaman-frekans uzayında temsil edilmiş ve farklı etiketleme protokolleri altında değerlendirilmiştir.

İlk deneylerde deprem öncesi kısa zaman pencerelerinin kronolojik testlerde ne ölçüde ayırt edilebildiği incelenmiştir. Bu protokolde pozitif sınıfın doğrulama ve test dönemlerinde oldukça seyrek kalması nedeniyle F1 ve precision metriklerinin sınırlı olduğu görülmüştür. Daha sonra temiz normal/anomali ayrımı protokolü kurulmuştur. Bu protokolde Mw $\geq$ 4,0 olaylarının $\pm$60 dakika çevresi anomali sınıfı, aynı büyüklükteki olaylardan en az 48 saat uzak pencereler ise temiz normal sınıf olarak tanımlanmıştır.

STFT temsillerinden çıkarılan 51 boyutlu özet öznitelikler ile HistGradientBoosting, ExtraTrees ve RandomForest modelleri karşılaştırılmıştır. En başarılı model olan HistGradientBoosting, dengeli test kümesinde 0,8124 accuracy, 0,8066 precision, 0,8219 recall, 0,8142 F1, 0,8979 ROC-AUC ve 0,9090 PR-AUC değerlerine ulaşmıştır. Elde edilen sonuçlar, deprem olayına yakın sismik pencereler ile olaylardan uzak temiz normal pencereler arasında öğrenilebilir zaman-frekans farklılıkları bulunduğunu göstermektedir. Çalışma deterministik deprem tahmini iddiası taşımamakta; sismik sinyallerde normal/anomali ayrımı için araştırma amaçlı bir prototip sunmaktadır.

\vspace{0.8cm}
\textbf{Anahtar Kelimeler:}

\begin{itemize}[leftmargin=*, itemsep=2pt]
    \item \textbf{Sismik Anomali Tespiti:} Sürekli sismik sinyallerdeki olağandışı davranışların otomatik olarak belirlenmesi.
    \item \textbf{STFT:} Sismik sinyallerin zaman-frekans içeriğini temsil etmek için kullanılan kısa süreli Fourier dönüşümü.
    \item \textbf{Zaman Serisi Analizi:} Zamana bağlı verilerin istatistiksel ve algoritmik yöntemlerle incelenmesi.
    \item \textbf{Makine Öğrenmesi:} Veriden çıkarılan özniteliklerle sınıflandırma ve skor üretme yöntemleri.
    \item \textbf{HistGradientBoosting:} Ağaç tabanlı gradyan artırma yöntemiyle çalışan sınıflandırma modeli.
    \item \textbf{Olay Çevresi Analizi:} Deprem olaylarının belirli zaman çevresindeki sismik pencerelerin incelenmesi.
\end{itemize}

% İçindekiler - Özet'ten hemen sonra
\tableofcontents

% -------------------- ŞEKİLLER DİZİNİ --------------------
\chapter*{ŞEKİLLER DİZİNİ}
\addcontentsline{toc}{chapter}{ŞEKİLLER DİZİNİ}

\begin{longtable}{p{2cm}p{11cm}r}
\textbf{Şekil} & \textbf{Açıklama} & \textbf{Sayfa} \\ \hline
\hyperref[fig:sistem_mimarisi]{Şekil 1.1} & DeepSeis Sistem Mimarisi Genel Görünümü & \pageref{fig:sistem_mimarisi} \\
\hyperref[fig:sismik_dalgalar]{Şekil 2.1} & Sismik Dalga Tipleri ve Yayılım Karakteristikleri & \pageref{fig:sismik_dalgalar} \\
\hyperref[fig:ml_anomaly_pipeline]{Şekil 2.2} & STFT Öznitelikleriyle Sınıflandırma Tabanlı Anomali Tespiti Yaklaşımı & \pageref{fig:ml_anomaly_pipeline} \\
\hyperref[fig:stft_ornek]{Şekil 3.1} & STFT Tabanlı Zaman-Frekans Dönüşümü Örneği & \pageref{fig:stft_ornek} \\
\hyperref[fig:clean_model_comparison]{Şekil 4.1} & Temiz Normal/Anomali Ayrımı Model Karşılaştırması & \pageref{fig:clean_model_comparison} \\
\hyperref[fig:clean_confusion]{Şekil 4.2} & HistGradientBoosting Modeli Karışıklık Matrisi & \pageref{fig:clean_confusion} \\
\hyperref[fig:clean_roc]{Şekil 4.3} & HistGradientBoosting Modeli ROC Eğrisi & \pageref{fig:clean_roc} \\
\hyperref[fig:clean_pr]{Şekil 4.4} & HistGradientBoosting Modeli Precision-Recall Eğrisi & \pageref{fig:clean_pr} \\
\end{longtable}

% -------------------- TABLOLAR DİZİNİ --------------------
\chapter*{TABLOLAR DİZİNİ}
\addcontentsline{toc}{chapter}{TABLOLAR DİZİNİ}

\begin{longtable}{p{2cm}p{11cm}r}
\textbf{Tablo} & \textbf{Açıklama} & \textbf{Sayfa} \\ \hline
\hyperref[tab:veri_ozellikleri]{Tablo 3.1} & Kullanılan Veri Kümesi Özeti & \pageref{tab:veri_ozellikleri} \\
\hyperref[tab:hyperparams]{Tablo 3.2} & Nihai Feature Model Ayarları & \pageref{tab:hyperparams} \\
\hyperref[tab:kronolojik_bolumleme]{Tablo 4.1} & Kronolojik Veri Bölümleme Özeti & \pageref{tab:kronolojik_bolumleme} \\
\hyperref[tab:kronolojik_sonuclar]{Tablo 4.2} & Kronolojik Testlerde Elde Edilen Temel Sonuçlar & \pageref{tab:kronolojik_sonuclar} \\
\hyperref[tab:temiz_deney_ozet]{Tablo 4.3} & Temiz Dengeli Deney Veri Özeti & \pageref{tab:temiz_deney_ozet} \\
\hyperref[tab:temiz_model_performans]{Tablo 4.4} & Temiz Normal/Anomali Ayrımı Model Performansı & \pageref{tab:temiz_model_performans} \\
\end{longtable}

% -------------------- SEMBOLLER VE KISALTMALAR --------------------
\chapter*{SEMBOLLER VE KISALTMALAR}
\addcontentsline{toc}{chapter}{SEMBOLLER VE KISALTMALAR}

% Satır aralığını azaltarak tek sayfaya sığdır
\renewcommand{\arraystretch}{0.95}

\begin{longtable}{@{}p{2.5cm}p{10.5cm}@{}}
\toprule
\textbf{Kısaltma} & \textbf{Açıklama} \\
\midrule
\endfirsthead
\toprule
\textbf{Kısaltma} & \textbf{Açıklama} \\
\midrule
\endhead
\endfoot
\bottomrule
\endlastfoot
AI & Artificial Intelligence (Yapay Zekâ) \\
ML & Machine Learning (Makine Öğrenmesi) \\
DL & Deep Learning (Derin Öğrenme) \\
TST & Time Series Transformer \\
TFT & Temporal Fusion Transformer \\
STFT & Short-Time Fourier Transform (Kısa Zamanlı Fourier Dönüşümü) \\
CWT & Continuous Wavelet Transform (Sürekli Dalgacık Dönüşümü) \\
FFT & Fast Fourier Transform (Hızlı Fourier Dönüşümü) \\
MSE & Mean Squared Error (Ortalama Kare Hatası) \\
MAE & Mean Absolute Error (Ortalama Mutlak Hata) \\
ROC-AUC & Receiver Operating Characteristic - Area Under Curve \\
PR-AUC & Precision-Recall - Area Under Curve \\
TP & True Positive (Doğru Pozitif) \\
FP & False Positive (Yanlış Pozitif) \\
FN & False Negative (Yanlış Negatif) \\
TN & True Negative (Doğru Negatif) \\
LSTM & Long Short-Term Memory \\
GRU & Gated Recurrent Unit \\
RNN & Recurrent Neural Network \\
CNN & Convolutional Neural Network \\
GPU & Graphics Processing Unit \\
STA/LTA & Short-Term Average / Long-Term Average \\
MiniSEED & Sismik Kayıt Veri Formatı \\
IRIS & Incorporated Research Institutions for Seismology \\
USGS & United States Geological Survey \\
AFAD & Afet ve Acil Durum Yönetimi Başkanlığı \\
Kandilli & Boğaziçi Üniversitesi Kandilli Rasathanesi \\
Mw & Moment Magnitüdü \\
$M_L$ & Lokal Magnitüd \\
ObsPy & Python Tabanlı Sismoloji Kütüphanesi \\
RMS & Root Mean Square \\
\end{longtable}

% Satır aralığını normale döndür
\renewcommand{\arraystretch}{1.0}


\newpage
\pagenumbering{arabic}
\setcounter{page}{1}

% -------------------- GİRİŞ --------------------
\chapter{GİRİŞ}

\section{Problem Tanımı}

Deprem, yüksek yıkıcılık potansiyeli nedeniyle hem insan hayatı hem de ekonomik açıdan kritik bir doğal afettir \cite{allen2009real}. Dünya genelinde her yıl ortalama 20.000'den fazla can kaybına neden olan depremler, aynı zamanda milyarlarca dolarlık ekonomik zarara yol açmaktadır. Türkiye, Kuzey Anadolu Fay Hattı ve Doğu Anadolu Fay Hattı gibi aktif tektonik yapılar üzerinde konumlanması nedeniyle tarih boyunca sayısız yıkıcı depreme maruz kalmıştır \cite{ambraseys2002seismicity}. 1999 Marmara Depremi (Mw 7.4) ve 2023 Kahramanmaraş Depremleri (Mw 7.7 ve Mw 7.6) bu yıkıcılığın en acı örnekleridir. Bu gerçeklik, deprem bilimi alanındaki araştırmaları ülkemiz için stratejik bir öncelik haline getirmektedir.

\subsection{Mevcut Erken Uyarı Sistemlerinin Sınırlamaları}

Günümüzde kullanılan erken uyarı sistemleri genellikle depremin ilk saniyelerinde P-dalgasını (birincil dalga) algılayarak, daha yıkıcı olan S-dalgasından (ikincil dalga) önce uyarı vermeye odaklanmaktadır \cite{kanamori2005real}. Bu sistemlerin çalışma prensibi şöyle özetlenebilir:

\begin{enumerate}
    \item \textbf{P-Dalgası Algılama:} Deprem kaynağından yayılan ilk sismik dalga olan P-dalgası, hızı nedeniyle istasyonlara ilk ulaşan dalgadır (yaklaşık 6-8 km/s).
    \item \textbf{Hızlı Analiz:} Sistem, P-dalgasının özelliklerinden (genlik, frekans içeriği) depremin büyüklüğünü ve konumunu tahmin eder.
    \item \textbf{Uyarı Yayını:} S-dalgası (yaklaşık 3-5 km/s) ve yüzey dalgaları varmadan önce saniyeler içinde uyarı verilir.
\end{enumerate}

Bu yaklaşım deprem başladıktan sonra birkaç saniyelik (genellikle 10-60 saniye) kazanım sağlamaktadır. Ancak bazı kritik sınırlamaları bulunmaktadır:

\begin{itemize}
    \item \textbf{Reaktif Yapı:} Sistem ancak deprem başladıktan sonra devreye girmektedir.
    \item \textbf{Kör Bölge:} Episantra yakın bölgelerde uyarı süresi çok kısa veya sıfır olabilir.
    \item \textbf{Tahmin Değil, Algılama:} Deprem öncesi herhangi bir öngörü sağlamamaktadır.
\end{itemize}

\subsection{Deprem Öncesi Anomali Tespitinin Önemi}

Deprem olmadan önceki saatler veya günler içinde yer kabuğundaki mikrosismik davranıştaki değişimleri modelleyerek ``olasılık artışı'' hakkında bir ön sinyal üretmek hâlâ bilimsel olarak zorlu ve tartışmalı bir konudur \cite{geller1997earthquakes}. Ancak son yıllarda yapılan araştırmalar, bazı depremlerin öncesinde şu tür anomalilerin gözlemlenebileceğini ortaya koymuştur:

\begin{itemize}
    \item \textbf{Mikrosismik Aktivite Değişimi:} Fay hattı çevresinde küçük depremlerin sıklığında artış veya azalma.
    \item \textbf{Sismik Hız Değişimi:} Yeraltı kayaç özelliklerinin stres altında değişmesi.
    \item \textbf{b-Değeri Anomalileri:} Gutenberg-Richter yasasındaki b parametresinin değişimi.
    \item \textbf{Frekans İçeriği Değişimi:} Sismik gürültünün spektral özelliklerinde kayma.
\end{itemize}

Bu tez çalışması, tamamen deterministik bir ``yarın deprem olacak'' iddiasında bulunmak yerine, sismik sinyallerdeki normal ve anomali örüntülerini otomatik olarak ayırt eden bir araştırma prototipi geliştirmeyi hedeflemektedir. Çalışmada özellikle deprem olayına yakın sismik pencereler ile olaylardan uzak temiz normal pencereler arasındaki zaman-frekans farkları incelenmektedir.

\begin{figure}[htbp]
    \centering
    \fitfig[max width=0.82\textwidth,max height=0.52\textheight]{figures/deepseis_system_architecture.png}
    \caption{DeepSeis Sistem Mimarisi Genel Görünümü}
    \label{fig:sistem_mimarisi}
\end{figure}

\section{Araştırma Soruları}

Bu çalışma aşağıdaki temel araştırma sorularını cevaplamayı amaçlamaktadır:

\begin{enumerate}
    \item \textbf{Anomali Varlığı:} Deprem öncesi belirli frekans bantlarında veya genlik değişimlerinde olağan dışı davranışlar (anomali) oluşmakta mıdır?
    
    \item \textbf{Model Etkinliği:} STFT tabanlı öznitelikler ve makine öğrenmesi modelleri, temiz normal ve olay çevresi anomali pencerelerini güvenilir şekilde ayırt edebilir mi?
    
    \item \textbf{İstatistiksel Anlamlılık:} Tespit edilen anomaliler ile gerçekleşen depremler arasında istatistiksel olarak anlamlı bir korelasyon var mıdır?
    
    \item \textbf{Öncü Süre:} Anomaliler tespit edilebiliyorsa, depremden ne kadar süre önce oluşmaktadır?
    
    \item \textbf{Genelleştirilebilirlik:} Geliştirilen model farklı bölgelere ve farklı büyüklükteki depremlere uygulanabilir mi?
\end{enumerate}

\section{Amaç ve Bilimsel Katkı}

Bu tezin temel amacı, sismik titreşim verilerini STFT tabanlı zaman-frekans temsillerine dönüştürerek:

\begin{itemize}
    \item Sürekli kaydedilen ivme/titreşim sinyallerinden anlamlı zaman-frekans öznitelikleri çıkarmak,
    \item Olay çevresi anomali pencereleri ile temiz normal pencereleri karşılaştırmak,
    \item Makine öğrenmesi ve derin öğrenme modellerinin bu görevdeki performansını ölçmek,
    \item En başarılı modeli çalışan bir web arayüzü içinde anomali skoru üretecek şekilde kullanmaktır.
\end{itemize}

\subsection{Neden STFT ve Makine Öğrenmesi Yaklaşımı?}

Bu çalışmada farklı derin öğrenme modelleri denenmiş, ancak nihai başarılı prototip STFT tabanlı öznitelikler ve ağaç tabanlı makine öğrenmesi modelleri üzerine kurulmuştur. Bu tercihin temel nedenleri şunlardır:

\begin{enumerate}
    \item \textbf{Zaman-frekans açıklığı:} STFT, sismik sinyalin hem zamansal hem de frekans içeriğini birlikte temsil eder.
    
    \item \textbf{Hızlı deney döngüsü:} Özet öznitelikler, farklı etiketleme protokollerinin hızlı biçimde karşılaştırılmasını sağlar.
    
    \item \textbf{Sınıf dengesizliğine dayanıklılık:} HistGradientBoosting, RandomForest ve ExtraTrees gibi modeller sınırlı ve dengesiz veri koşullarında güçlü referans sonuçlar üretebilir.
    
    \item \textbf{Uygulama kolaylığı:} Eğitilen model küçük boyutlu \texttt{joblib} çıktısı olarak saklanabilir ve web arayüzünde hızlı çalıştırılabilir.
\end{enumerate}

\subsection{Bilimsel Katkılar}

Bu çalışmanın literatüre sağlayacağı başlıca bilimsel katkılar şunlardır:

\begin{enumerate}
    \item \textbf{Model Karşılaştırması:} Derin öğrenme modelleri ile STFT özniteliklerine dayalı makine öğrenmesi modellerinin sismik anomali tespitindeki performanslarının karşılaştırılması.
    
    \item \textbf{Çoklu Girdi Temsilleri:} Sismik verinin ham dalga formundan (raw waveform) spektral temsillere (spektrogram, STFT, Mel benzeri bant güçleri) dönüştürülüp farklı girdi biçimlerinin denenmesi ve karşılaştırılması.
    
    \item \textbf{Nicel Değerlendirme:} Normal/anomali ayrımının accuracy, precision, recall, F1, ROC-AUC ve PR-AUC metrikleriyle ölçülmesi.
    
    \item \textbf{Yerel Uygulama:} Türkiye odaklı gerçek istasyon verileri kullanılarak yerel bir erken uyarı / risk artışı analiz prototipinin hazırlanması.
    
    \item \textbf{Çalışan Prototip:} Geliştirilen modelin Flask tabanlı arayüzde anomali skoru üretecek biçimde çalıştırılması.
\end{enumerate}

\section{Çalışmanın Kapsamı ve Sınırları}

Bu çalışma aşağıdaki kapsam dahilinde yürütülmektedir:

\begin{itemize}
    \item \textbf{Coğrafi Kapsam:} Türkiye'deki seçili sismik istasyonlar (Doğu Anadolu Fay Hattı çevresi)
    \item \textbf{Zaman Kapsamı:} 6-12 aylık sürekli sismik kayıtlar
    \item \textbf{Büyüklük Kapsamı:} Ana deneylerde Mw $\geq$ 4.0 büyüklüğündeki depremler
\end{itemize}

Çalışmanın bilinen sınırlamaları:

\begin{itemize}
    \item Bu çalışma, deterministik deprem tahmini iddiasında bulunmamaktadır.
    \item Sonuçlar belirli bir bölge ve zaman dilimi için geçerlidir; genelleştirme ayrı çalışma gerektirir.
    \item Tespit edilen anomalilerin tamamı tektonik kökenli olmayabilir.
\end{itemize}

\section{Tezin Organizasyonu}

Bu tez altı bölümden oluşmaktadır:

\textbf{Bölüm 1 (Giriş):} Problem tanımı, mevcut sistemlerin sınırlamaları, araştırma soruları ve çalışmanın amaçları sunulmaktadır.

\textbf{Bölüm 2 (Literatür Taraması):} Sismik veri analizi, anomali tespiti yöntemleri, makine öğrenmesi ve derin öğrenme yaklaşımları hakkındaki mevcut çalışmalar kapsamlı şekilde incelenmektedir.

\textbf{Bölüm 3 (Yöntem):} Veri kaynakları, veri ön işleme teknikleri, model mimarisi, eğitim stratejisi ve değerlendirme metrikleri detaylı olarak açıklanmaktadır.

\textbf{Bölüm 4 (Bulgular ve Tartışma):} Deneysel sonuçlar sunulmakta, model performansları karşılaştırılmakta ve bulgular literatür ışığında değerlendirilmektedir.

\textbf{Bölüm 5 (Sonuç ve Öneriler):} Çalışmanın sonuçları özetlenmekte, sınırlılıklar tartışılmakta ve gelecek çalışmalar için öneriler sunulmaktadır.

% -------------------- LİTERATÜR TARAMASI --------------------
\chapter{LİTERATÜR TARAMASI}

\section{Sismik Veri Analizi ve Deprem Erken Uyarı Sistemleri}

Sismik veri analizi, yer kabuğundaki titreşimlerin kaydedilmesi, işlenmesi ve yorumlanmasını kapsayan multidisipliner bir alandır. Deprem dalgaları, P-dalgaları (birincil, boyuna dalgalar) ve S-dalgaları (ikincil, enine dalgalar) olmak üzere iki ana kategoriye ayrılmaktadır \cite{stein2009introduction}. P-dalgaları daha hızlı yayılırken, S-dalgaları daha fazla enerji taşımakta ve yapısal hasara neden olmaktadır.

\begin{figure}[htbp]
    \centering
    \fitfig[max width=0.82\textwidth,max height=0.48\textheight]{figures/seismic_wave_types.png}
    \caption{Sismik Dalga Tipleri ve Yayılım Karakteristikleri}
    \label{fig:sismik_dalgalar}
\end{figure}

Geleneksel erken uyarı sistemleri, P-dalgasının algılanmasından S-dalgasının varışına kadar geçen süreyi kullanarak uyarı vermektedir \cite{allen2009real}. Japonya'daki JMA (Japan Meteorological Agency) sistemi ve ABD'deki ShakeAlert bu yaklaşımın önde gelen örnekleridir \cite{given2014technical}. Türkiye'de ise İstanbul için geliştirilen erken uyarı sistemi benzer prensiplerle çalışmaktadır \cite{erdik2003earthquake}.

\subsection{Klasik Sinyal İşleme Yöntemleri}

Sismik sinyallerde olay tespiti için yaygın olarak kullanılan klasik yöntemler şunlardır:

\begin{itemize}
    \item \textbf{STA/LTA (Short-Term Average / Long-Term Average):} Kısa vadeli ve uzun vadeli sinyal ortalamalarının oranını hesaplayarak ani değişimleri tespit eder \cite{allen1978automatic}. Bu yöntem basit ve hesaplama açısından verimli olmasına rağmen, gürültülü ortamlarda yüksek yanlış alarm oranına sahiptir.
    
    \item \textbf{Eşik Tabanlı Tetikleyiciler:} Sinyal genliği belirli bir eşik değerini aştığında olay kaydı başlatılır. Bu yöntem, eşik değerinin seçimine oldukça duyarlıdır.
    
    \item \textbf{Spektral Analiz:} Fourier dönüşümü kullanılarak sinyalin frekans bileşenleri analiz edilir \cite{oppenheim1999discrete}. Deprem sinyalleri karakteristik frekans örüntüleri sergilemektedir.
\end{itemize}

\section{Makine Öğrenmesi ve Derin Öğrenme Yaklaşımları}

Son yıllarda makine öğrenmesi ve derin öğrenme yöntemleri, sismik veri analizinde önemli başarılar elde etmiştir \cite{kong2019machine}. Bu yaklaşımlar, geleneksel yöntemlerin sınırlamalarını aşma potansiyeli taşımaktadır.

\subsection{Tekrarlayan Sinir Ağları (RNN)}

LSTM (Long Short-Term Memory) ve GRU (Gated Recurrent Unit) mimarileri, zaman serisi verilerindeki sekansiyel bağımlılıkları modellemek için tasarlanmıştır \cite{hochreiter1997long}. Sismik sinyallerin analizi için yaygın olarak kullanılmakta olup, özellikle P-dalgası tespiti ve faz ayrıştırma görevlerinde başarılı sonuçlar elde edilmiştir \cite{ross2018generalized}.

Ancak RNN tabanlı modeller, uzun sekanslardan kaynaklanan gradyan kaybolması problemi nedeniyle uzun vadeli bağımlılıkları modellemede sınırlı kalmaktadır \cite{bengio1994learning}. Sismik sinyallerde deprem öncesi anomalilerin saatler hatta günler öncesinde oluşabileceği düşünüldüğünde, bu sınırlama kritik bir öneme sahiptir.

\subsection{Evrişimli Sinir Ağları (CNN)}

CNN mimarileri, spektrogram ve zaman-frekans temsillerinin analizi için etkili olduğu gösterilmiştir \cite{perol2018convolutional}. ConvNetQuake ve benzeri modeller, sismik olayların tespiti ve sınıflandırılmasında yüksek doğruluk oranları elde etmiştir.

\section{Transformer Mimarisi ve Zaman Serisi Analizi}

Transformer mimarisi, Vaswani ve arkadaşları tarafından 2017 yılında önerilmiş olup, doğal dil işleme alanında devrim yaratmıştır \cite{vaswani2017attention}. Self-attention mekanizması sayesinde, sekans içindeki herhangi iki pozisyon arasındaki bağımlılıkları doğrudan modelleyebilmektedir.

\subsection{Attention Mekanizması}

Attention mekanizması, girdi sekansındaki her elemanın diğer elemanlarla olan ilişkisini öğrenmektedir. Matematiksel olarak:

\begin{equation}
\text{Attention}(Q, K, V) = \text{softmax}\left(\frac{QK^T}{\sqrt{d_k}}\right)V
\end{equation}

Burada $Q$ (Query), $K$ (Key) ve $V$ (Value) matrisleri girdi sekansından türetilmekte, $d_k$ ise anahtar vektörlerinin boyutunu temsil etmektedir.

\begin{figure}[htbp]
    \centering
    \fitfig[max width=0.82\textwidth,max height=0.48\textheight]{figures/ml_anomaly_pipeline.png}
    \caption{STFT Öznitelikleriyle Sınıflandırma Tabanlı Anomali Tespiti Yaklaşımı}
    \label{fig:ml_anomaly_pipeline}
\end{figure}

\subsection{Zaman Serisi için Transformer Varyantları}

Zaman serisi analizi için özelleştirilmiş Transformer varyantları geliştirilmiştir:

\begin{enumerate}
    \item \textbf{Time Series Transformer (TST):} Doğrudan zaman serisi verileri için tasarlanmış olup, pozisyonel kodlama ve maskeleme stratejileri içermektedir \cite{li2019enhancing}.
    
    \item \textbf{Informer:} Uzun sekans tahmini için ProbSparse self-attention mekanizması kullanan verimli bir Transformer varyantıdır \cite{zhou2021informer}. $O(L \log L)$ karmaşıklığı ile uzun sekanslarda etkili çalışmaktadır.
    
    \item \textbf{Temporal Fusion Transformer (TFT):} Çok değişkenli zaman serisi tahmini için tasarlanmış olup, statik kovaryantlar, geçmiş gözlemler ve gelecek bilinen girdileri birleştirmektedir \cite{lim2021temporal}.
    
    \item \textbf{Autoformer:} Auto-correlation mekanizması ile mevsimsel ve trend bileşenlerini ayrıştıran bir mimaridir \cite{wu2021autoformer}.
\end{enumerate}

\section{Anomali Tespiti Yaklaşımları}

Zaman serisi verilerinde anomali tespiti, çeşitli yaklaşımlarla gerçekleştirilebilmektedir:

\subsection{Yeniden Yapılandırma Tabanlı Yöntemler}

Autoencoder mimarileri, normal veriyi sıkıştırıp yeniden yapılandırmayı öğrenmektedir. Yeniden yapılandırma hatası yüksek olan örnekler anomali olarak işaretlenmektedir \cite{sakurada2014anomaly}:

\begin{equation}
\text{Anomali Skoru} = ||x - \hat{x}||^2
\end{equation}

Burada $x$ orijinal girdi, $\hat{x}$ ise yeniden yapılandırılmış çıktıdır.

\subsection{Tahmin Tabanlı Yöntemler}

Model, bir sonraki zaman adımını veya pencereyi tahmin etmek üzere eğitilmektedir. Gerçek değer ile tahmin arasındaki fark, anomali skorunu oluşturmaktadır.

\subsection{Sınıflandırma Tabanlı Yöntemler}

İkili sınıflandırma yaklaşımıyla, her zaman penceresi ``normal'' veya ``anormal'' olarak etiketlenmektedir. Bu yaklaşım, etiketli veriye ihtiyaç duymaktadır.

\section{Sismik Anomali ve Deprem Öncesi Sinyaller}

Deprem öncesi anomali tespiti, sismoloji alanında tartışmalı bir konu olmaya devam etmektedir \cite{geller1997earthquakes}. Bazı çalışmalar, deprem öncesi sismik aktivitede değişimler raporlamış olsa da, bu bulguların evrenselliği ve güvenilirliği sorgulanmaktadır.

\subsection{Mikrosisimik Aktivite Değişimleri}

Deprem öncesi dönemde mikrosismik aktivitede artış veya sessizlik dönemleri gözlemlenebilmektedir \cite{scholz1990mechanics}. Bu değişimler, fay hattındaki gerilim birikiminin göstergesi olabilir.

\subsection{Frekans İçeriği Değişimleri}

Bazı araştırmalar, deprem öncesi sismik sinyallerin frekans içeriğinde değişimler tespit etmiştir. Bu değişimler, kayaç özelliklerinin stres altında değişmesiyle ilişkilendirilebilir.

\section{Literatür Özeti}

\begin{table}[htbp]
\centering
\small
\caption{Literatürdeki İlgili Çalışmaların Karşılaştırması}
\begin{adjustbox}{center,max width=\textwidth}
\begin{tabularx}{0.96\textwidth}{p{2.45cm}p{2.45cm}p{3.0cm}X}
\toprule
\textbf{Çalışma} & \textbf{Yöntem} & \textbf{Veri Seti} & \textbf{Sonuçlar} \\
\midrule
Ross et al., 2018 & CNN + RNN & Southern California & P-dalgası tespitinde yüksek doğruluk \\
Perol et al., 2018 & ConvNetQuake & Oklahoma & Küçük depremlerin başarılı tespiti \\
Mousavi et al., 2020 & EQTransformer & Global & Çoklu görev öğrenmede başarı \\
Zhou et al., 2021 & Informer & Çeşitli zaman serileri & Uzun vadeli tahminde üstün performans \\
\bottomrule
\end{tabularx}
\end{adjustbox}
\label{tab:literatur}
\end{table}

Bu literatür taraması ışığında, sismik anomali analizinde hem derin öğrenme hem de öznitelik tabanlı makine öğrenmesi yaklaşımlarının önemli bir araştırma alanı oluşturduğu görülmektedir. Bu tez, STFT tabanlı zaman-frekans öznitelikleri ile temiz normal/anomali ayrımı yapabilen uçtan uca bir DeepSeis prototipi geliştirerek bu alana uygulamalı bir katkı sunmayı hedeflemektedir.

% -------------------- YÖNTEM --------------------
\chapter{YÖNTEM}

Bu bölümde DeepSeis sisteminin veri işleme, etiketleme, modelleme ve değerlendirme adımları açıklanmaktadır. Çalışma, sismik sinyalden zaman-frekans temsili üretme, farklı deney protokollerini karşılaştırma ve en başarılı modeli web arayüzünde çalıştırma aşamalarından oluşmaktadır.

\section{Veri Kaynakları}

\subsection{Sürekli Sismik İstasyon Verileri}

Çalışmada sürekli sismik istasyon kayıtları kullanılmıştır. Ham veriler MiniSEED formatında saklanmakta ve istasyon bileşenleri HHZ, HHN ve HHE kanalları üzerinden temsil edilmektedir. Bu tez kapsamında güçlü ve tutarlı ilk deney seti HHZ düşey bileşeni üzerinden oluşturulmuştur. Üç bileşenli modelleme için kanalların zaman damgası seviyesinde hizalanması gerektiğinden, çok bileşenli füzyon gelecek çalışma olarak bırakılmıştır.

\subsection{Deprem Kataloğu}

Etiket üretimi için aynı zaman aralığına karşılık gelen deprem kataloğu kullanılmıştır. Katalog kayıtları olay zamanı, büyüklük, konum ve derinlik bilgilerini içermektedir. Deneylerde özellikle Mw $\geq$ 4,0 olayları kullanılmıştır; çünkü daha küçük olaylar gürültüyle daha kolay karışabilmekte, daha büyük olaylar ise veri aralığında az sayıda bulunabilmektedir.

\section{Veri Ön İşleme}

\subsection{Pencereleme}

Sürekli sismik kayıtlar kısa zaman pencerelerine ayrılmıştır. Her pencere bağımsız bir örnek olarak ele alınmış ve daha sonra STFT tabanlı zaman-frekans temsiline dönüştürülmüştür. İşlenmiş HHZ veri kümesi toplam 162.337 pencere içermektedir.

\begin{table}[htbp]
\centering
\small
\caption{Kullanılan Veri Kümesi Özeti}
\begin{adjustbox}{center,max width=\textwidth}
\begin{tabular}{ll}
\toprule
\textbf{Özellik} & \textbf{Değer} \\
\midrule
Kanal & HHZ \\
Toplam pencere & 162.337 \\
STFT boyutu & 129 $\times$ 95 \\
Kronolojik eğitim örneği & 123.422 \\
Kronolojik doğrulama örneği & 17.077 \\
Kronolojik test örneği & 21.838 \\
\bottomrule
\end{tabular}
\end{adjustbox}
\label{tab:veri_ozellikleri}
\end{table}

\subsection{Spektral Dönüşüm}

Sismik pencereler kısa süreli Fourier dönüşümü (STFT) ile zaman-frekans uzayına taşınmıştır. STFT, sinyalin farklı zaman aralıklarında hangi frekans bileşenlerini içerdiğini gösterdiği için sismik anomali analizi açısından uygundur:

\begin{equation}
STFT\{x(t)\}(\tau, f) = \int_{-\infty}^{\infty} x(t) w(t-\tau) e^{-j2\pi ft} dt
\end{equation}

Burada $w(t)$ pencere fonksiyonu, $\tau$ zaman kayması ve $f$ frekans değişkenidir.

\begin{figure}[htbp]
    \centering
    \fitfig[max width=0.78\textwidth,max height=0.48\textheight]{figures/sample_spectrogram.png}
    \caption{STFT tabanlı zaman-frekans dönüşümü örneği}
    \label{fig:stft_ornek}
\end{figure}

\section{Öznitelik Çıkarımı}

Derin öğrenme modellerinde STFT matrisi doğrudan giriş olarak kullanılmıştır. Feature tabanlı modeller için ise her STFT penceresinden 51 boyutlu özet öznitelik vektörü çıkarılmıştır. Bu öznitelikler aşağıdaki gruplardan oluşmaktadır:

\begin{itemize}
    \item Spektrogram geneli için ortalama, standart sapma, minimum, maksimum ve yüzdelik değerler,
    \item Frekans profili ve zaman profili istatistikleri,
    \item Sekiz frekans bandı ve sekiz zaman bandı ortalamaları,
    \item Frekans ve zaman yönlü gradyan özetleri,
    \item 4$\times$4 zaman-frekans ızgarası üzerinde bölgesel enerji ortalamaları.
\end{itemize}

Bu yaklaşım, spektrogramdaki enerji dağılımını kompakt ve hızlı eğitilebilir bir temsil haline getirmiştir.

\section{Etiketleme Protokolleri}

Çalışmada üç farklı problem tanımı değerlendirilmiştir.

\subsection{Deprem Öncesi Kısa Pencere}

İlk deneylerde Mw $\geq$ 4,0 olaylarından önceki 1 saatlik pencereler pozitif sınıf olarak etiketlenmiştir:

\begin{equation}
y_i = \begin{cases}
1, & t_{eq} - 1\text{ saat} \leq t_i \leq t_{eq} \\
0, & \text{diğer}
\end{cases}
\end{equation}

Bu kurulum, deprem öncesi kısa dönemli sinyal aramayı hedeflediği için en zor deney protokollerinden biridir.

\subsection{Olay Çevresi Anomali Penceresi}

İkinci protokolde Mw $\geq$ 4,0 olaylarının $\pm$60 dakika çevresindeki pencereler pozitif sınıf olarak tanımlanmıştır:

\begin{equation}
y_i = \begin{cases}
1, & |t_i - t_{eq}| \leq 60\text{ dakika} \\
0, & \text{diğer}
\end{cases}
\end{equation}

Bu protokol, doğrudan deprem tahmini değil, olay çevresi sismik aktiviteyi ayırt etme görevidir.

\subsection{Temiz Normal / Anomali Ayrımı}

Nihai başarılı deneyde pozitif sınıf, Mw $\geq$ 4,0 olaylarının $\pm$60 dakika çevresinden; negatif sınıf ise herhangi bir Mw $\geq$ 4,0 olayından en az 48 saat uzak pencerelerden seçilmiştir:

\begin{equation}
y_i = \begin{cases}
1, & |t_i - t_{eq}| \leq 60\text{ dakika} \\
0, & |t_i - t_{eq}| \geq 48\text{ saat}
\end{cases}
\end{equation}

Bu protokol, modelin temiz normal dönemler ile olay çevresi anomali dönemleri arasındaki ayırt edilebilir örüntüleri öğrenip öğrenemediğini ölçmektedir.

\section{Modelleme Yaklaşımı}

\subsection{Derin Öğrenme Modelleri}

Proje kapsamında LSTM, GRU, CNN, Transformer ve hibrit CNN-Transformer mimarileri denenmiştir. CNN modeli STFT spektrogramını iki boyutlu görüntü gibi işlemiştir. Transformer tabanlı modeller ise zaman-frekans temsilinden uzun dönemli ilişkiler öğrenmeyi hedeflemiştir. Kronolojik testlerde sınıf dengesizliği ve dağılım kayması nedeniyle bu modellerin F1 skorları sınırlı kalmıştır.

\subsection{Feature Tabanlı Modeller}

STFT özet öznitelikleri üzerinde Logistic Regression, HistGradientBoosting, ExtraTrees ve RandomForest modelleri denenmiştir. En başarılı ve dengeli sonuç HistGradientBoosting modeliyle elde edilmiştir. Model, her pencere için anomali olasılığı üretmekte; karar eşiği doğrulama kümesinde F1 skorunu en iyi yapan değere göre seçilmektedir.

\begin{table}[htbp]
\centering
\small
\caption{Nihai Feature Model Ayarları}
\begin{adjustbox}{center,max width=\textwidth}
\begin{tabularx}{0.88\textwidth}{lX}
\toprule
\textbf{Parametre} & \textbf{Değer} \\
\midrule
Model & HistGradientBoostingClassifier \\
Öznitelik sayısı & 51 \\
Pozitif sınıf & Mw $\geq$ 4,0 olayının $\pm$60 dakikası \\
Negatif sınıf & Olaylardan en az 48 saat uzak pencereler \\
Test örneği & 7.826 \\
Karar eşiği & 0,4433 \\
\bottomrule
\end{tabularx}
\end{adjustbox}
\label{tab:hyperparams}
\end{table}

\section{Değerlendirme Metrikleri}

Model performansı accuracy, precision, recall, F1, ROC-AUC ve PR-AUC metrikleriyle değerlendirilmiştir:

\begin{align}
\text{Precision} &= \frac{TP}{TP + FP} \\
\text{Recall} &= \frac{TP}{TP + FN} \\
\text{F1-Score} &= 2 \times \frac{\text{Precision} \times \text{Recall}}{\text{Precision} + \text{Recall}}
\end{align}

Sınıf dengesizliği bulunduğunda accuracy tek başına yeterli değildir. Bu nedenle özellikle F1, recall, precision ve PR-AUC değerleri birlikte yorumlanmıştır.

\section{Uygulama Arayüzü}

Eğitilen en başarılı model \texttt{joblib} formatında kaydedilmiş ve Flask tabanlı DeepSeis arayüzüne bağlanmıştır. Arayüz, test pencerelerini sırayla okuyarak her pencere için anomali olasılığı üretmekte ve karar eşiği aşıldığında görsel uyarı vermektedir. Böylece model yalnızca raporlanan bir deney sonucu olarak kalmamış, çalışan bir prototip sistem içinde kullanılabilir hale getirilmiştir.

\section{Kullanılan Teknolojiler}

\begin{itemize}
    \item \textbf{Python:} Veri işleme, model eğitimi ve deney yönetimi
    \item \textbf{NumPy / SciPy:} Sayısal hesaplama ve sinyal işleme
    \item \textbf{PyTorch:} Derin öğrenme modelleri
    \item \textbf{Scikit-learn:} Feature tabanlı sınıflandırıcılar ve metrikler
    \item \textbf{Matplotlib:} Tez görselleri
    \item \textbf{Flask:} Web arayüzü ve model servisi
\end{itemize}

% -------------------- BULGULAR VE TARTIŞMA --------------------
\chapter{BULGULAR VE TARTIŞMA}

Bu bölümde DeepSeis projesi kapsamında yürütülen deneylerin sonuçları sunulmaktadır. Deneyler iki ayrı değerlendirme düzeyinde ele alınmıştır. İlk düzeyde, veri zaman sırası korunarak deprem öncesi veya olay çevresi pencerelerinin gelecekteki dönemlere genellenip genellenemediği incelenmiştir. İkinci düzeyde ise temiz normal ve olay çevresi anomali pencereleri dengeli biçimde ayrılarak modelin sismik sinyaldeki belirgin normal/anomali örüntülerini ayırt etme kapasitesi değerlendirilmiştir.

\section{Deneysel Kurulum}

\subsection{Veri Seti İstatistikleri}

Çalışmada HHZ düşey bileşeninden elde edilen STFT pencereleri kullanılmıştır. İşlenmiş veri kümesi toplam 162.337 pencere içermektedir. Kronolojik bölümleme sonucunda eğitim, doğrulama ve test kümeleri sırasıyla 123.422, 17.077 ve 21.838 örnekten oluşmuştur.

\begin{table}[htbp]
\centering
\small
\caption{Kronolojik Veri Bölümleme Özeti}
\begin{adjustbox}{center,max width=\textwidth}
\begin{tabular}{lrrrr}
\toprule
\textbf{Bölüm} & \textbf{Pencere} & \textbf{Pozitif} & \textbf{Negatif} & \textbf{Pozitif Oranı} \\
\midrule
Eğitim & 123.422 & 11.171 & 112.251 & \%9,05 \\
Doğrulama & 17.077 & 179 & 16.898 & \%1,05 \\
Test & 21.838 & 240 & 21.598 & \%1,10 \\
\bottomrule
\end{tabular}
\end{adjustbox}
\label{tab:kronolojik_bolumleme}
\end{table}

Tablo \ref{tab:kronolojik_bolumleme}, zaman sıralı bölümlemede pozitif sınıf oranının eğitim döneminden doğrulama ve test dönemlerine belirgin biçimde düştüğünü göstermektedir. Bu durum, deprem olaylarının veri zaman aralığı boyunca homojen dağılmadığını ve problemin yalnızca sınıf dengesizliği değil, aynı zamanda dağılım kayması içerdiğini göstermektedir.

\subsection{Etiketleme Protokolleri}

İlk deneylerde Mw $\geq$ 4,0 olaylarından önceki 1 saatlik pencereler pozitif sınıf olarak işaretlenmiştir. Bu kurulum, deprem öncesi kısa dönemli sinyal arama açısından daha iddialı ve zor bir problemdir. Daha sonra olay çevresi anomali tespiti için Mw $\geq$ 4,0 olaylarının $\pm$60 dakika çevresi pozitif sınıf olarak tanımlanmıştır. Son başarılı deneyde negatif sınıf, herhangi bir Mw $\geq$ 4,0 olayından en az 48 saat uzak pencerelerden seçilmiştir. Böylece modelin belirsiz geçiş bölgeleri yerine temiz normal ve temiz olay çevresi örüntülerini ayırt etmesi hedeflenmiştir.

\section{Kronolojik Test Sonuçları}

Kronolojik test protokolü, modelin geçmiş dönemlerde öğrenilen örüntüleri daha sonraki zaman aralıklarına aktarabilmesini ölçer. Bu protokol veri sızıntısını azaltması bakımından önemlidir; ancak olayların zamansal yoğunluğu değiştiğinde performans metrikleri sert biçimde düşebilmektedir.

\begin{table}[htbp]
\centering
\small
\caption{Kronolojik Testlerde Elde Edilen Temel Sonuçlar}
\begin{adjustbox}{center,max width=\textwidth}
\begin{tabularx}{0.96\textwidth}{Xcccccc}
\toprule
\textbf{Model / Protokol} & \textbf{Accuracy} & \textbf{Precision} & \textbf{Recall} & \textbf{F1} & \textbf{ROC-AUC} & \textbf{PR-AUC} \\
\midrule
CNN, Mw$\geq$4,0, 1 saat önce & 0,9015 & 0,0286 & 0,2417 & 0,0511 & 0,6324 & 0,0286 \\
CNN Focal Loss, Mw$\geq$4,0, 1 saat önce & 0,9764 & 0,0521 & 0,0667 & 0,0585 & 0,5538 & 0,0294 \\
Lojistik Regresyon, STFT öznitelikleri & 0,8177 & 0,0205 & 0,3333 & 0,0386 & 0,6484 & 0,0195 \\
\bottomrule
\end{tabularx}
\end{adjustbox}
\label{tab:kronolojik_sonuclar}
\end{table}

Tablo \ref{tab:kronolojik_sonuclar} sonuçları, kronolojik test senaryosunda accuracy değerinin tek başına yanıltıcı olabileceğini göstermektedir. Test kümesinde pozitif sınıf oranı yaklaşık \%1 olduğu için negatif sınıfı baskın tahmin eden modeller yüksek accuracy üretebilmekte, ancak precision ve F1 değerleri düşük kalmaktadır. Bu nedenle bu çalışmada ROC-AUC, PR-AUC, recall, precision ve karışıklık matrisi birlikte raporlanmıştır.

\section{Temiz Normal / Anomali Ayrımı}

\subsection{Dengeli Temiz Test Protokolü}

Kronolojik protokolde gözlenen dağılım kaymasını ayrıca incelemek için ikinci bir deney protokolü kurulmuştur. Bu protokolde pozitif sınıf Mw $\geq$ 4,0 deprem olaylarının $\pm$60 dakika çevresindeki pencerelerden, negatif sınıf ise herhangi bir Mw $\geq$ 4,0 olayından en az 48 saat uzak pencerelerden oluşturulmuştur. Pozitif ve negatif örnekler dengeli seçilmiş; eğitim, doğrulama ve test ayrımı stratified random split ile yapılmıştır.

\begin{table}[htbp]
\centering
\small
\caption{Temiz Dengeli Deney Veri Özeti}
\begin{adjustbox}{center,max width=\textwidth}
\begin{tabularx}{0.88\textwidth}{Xr}
\toprule
\textbf{Özellik} & \textbf{Değer} \\
\midrule
Pozitif tanımı & Mw $\geq$ 4,0 olayının $\pm$60 dakikası \\
Negatif tanımı & Mw $\geq$ 4,0 olaylarından en az 48 saat uzak \\
Toplam pozitif pencere & 19.564 \\
Temiz negatif havuzu & 31.818 \\
Eğitim örneği & 25.041 \\
Doğrulama örneği & 6.261 \\
Test örneği & 7.826 \\
Test pozitif / negatif & 3.913 / 3.913 \\
\bottomrule
\end{tabularx}
\end{adjustbox}
\label{tab:temiz_deney_ozet}
\end{table}

Bu protokol doğrudan deprem tahmini iddiası taşımaz. Amaç, deprem olayına yakın sismik aktivite pencereleri ile olaylardan uzak normal referans pencereleri arasında ayırt edilebilir bir zaman-frekans örüntüsü bulunup bulunmadığını ölçmektir.

\subsection{Model Performans Karşılaştırması}

STFT spektrogramlarından 51 boyutlu özet öznitelik çıkarılmıştır. Bu öznitelikler küresel enerji istatistikleri, frekans bant ortalamaları, zaman bant ortalamaları, gradyan istatistikleri ve kaba zaman-frekans ızgara ortalamalarından oluşmaktadır. HistGradientBoosting, ExtraTrees ve RandomForest modelleri aynı protokolde karşılaştırılmıştır.

\begin{table}[htbp]
\centering
\small
\caption{Temiz Normal/Anomali Ayrımı Model Performansı}
\begin{adjustbox}{center,max width=\textwidth}
\begin{tabularx}{0.96\textwidth}{Xcccccc}
\toprule
\textbf{Model} & \textbf{Accuracy} & \textbf{Precision} & \textbf{Recall} & \textbf{F1} & \textbf{ROC-AUC} & \textbf{PR-AUC} \\
\midrule
HistGradientBoosting & \textbf{0,8124} & \textbf{0,8066} & 0,8219 & \textbf{0,8142} & \textbf{0,8979} & \textbf{0,9090} \\
ExtraTrees & 0,8018 & 0,7799 & \textbf{0,8410} & 0,8093 & 0,8876 & 0,8965 \\
RandomForest & 0,8117 & 0,8034 & 0,8252 & 0,8142 & 0,8935 & 0,9044 \\
\bottomrule
\end{tabularx}
\end{adjustbox}
\label{tab:temiz_model_performans}
\end{table}

Tablo \ref{tab:temiz_model_performans} incelendiğinde en dengeli sonucun HistGradientBoosting modeli tarafından üretildiği görülmektedir. Bu model test kümesinde 0,8124 accuracy, 0,8066 precision, 0,8219 recall ve 0,8142 F1 skoruna ulaşmıştır. ROC-AUC değerinin 0,8979 ve PR-AUC değerinin 0,9090 olması, modelin yalnızca tek bir eşikte değil, skor sıralaması bakımından da güçlü ayrım yaptığını göstermektedir.

\begin{figure}[htbp]
    \centering
    \fitfig[max width=0.88\textwidth,max height=0.50\textheight]{figures/clean_model_comparison.png}
    \caption{Temiz normal/anomali ayrımı için model karşılaştırması}
    \label{fig:clean_model_comparison}
\end{figure}

\subsection{Karışıklık Matrisi}

En iyi model olarak seçilen HistGradientBoosting sınıflandırıcısının test kümesi karışıklık matrisi Şekil \ref{fig:clean_confusion}'da verilmiştir.

\begin{figure}[htbp]
    \centering
    \fitfig[max width=0.54\textwidth,max height=0.42\textheight]{figures/clean_hist_gradient_confusion_matrix.png}
    \caption{HistGradientBoosting modeli karışıklık matrisi}
    \label{fig:clean_confusion}
\end{figure}

Model, 3.913 pozitif test örneğinin 3.216'sını doğru biçimde anomali olarak sınıflandırmıştır. Normal sınıfta ise 3.913 örneğin 3.142'si doğru sınıflandırılmıştır. Yanlış negatif sayısı 697, yanlış pozitif sayısı 771 olarak ölçülmüştür. Bu sonuç, modelin recall tarafında güçlü olduğunu ve olay çevresi pencerelerinin büyük bölümünü yakaladığını göstermektedir.

\subsection{ROC ve Precision-Recall Analizi}

Şekil \ref{fig:clean_roc} ve Şekil \ref{fig:clean_pr}, HistGradientBoosting modelinin skor tabanlı ayrım gücünü göstermektedir.

\begin{figure}[htbp]
    \centering
    \fitfig[max width=0.58\textwidth,max height=0.42\textheight]{figures/clean_hist_gradient_roc_curve.png}
    \caption{HistGradientBoosting modeli ROC eğrisi}
    \label{fig:clean_roc}
\end{figure}

\begin{figure}[htbp]
    \centering
    \fitfig[max width=0.58\textwidth,max height=0.42\textheight]{figures/clean_hist_gradient_pr_curve.png}
    \caption{HistGradientBoosting modeli precision-recall eğrisi}
    \label{fig:clean_pr}
\end{figure}

ROC-AUC değerinin 0,8979 olması, modelin normal ve olay çevresi anomali pencerelerini yüksek oranda ayırabildiğini göstermektedir. PR-AUC değerinin 0,9090 olması ise dengeli test kümesinde pozitif sınıfın sıralama kalitesinin yüksek olduğunu ortaya koymaktadır.

\section{Tartışma}

\subsection{Ana Bulgular}

Deney sonuçları üç temel bulguya işaret etmektedir:

\begin{enumerate}
    \item \textbf{Kronolojik tahmin senaryosu zordur:} Veri zaman sırası korunduğunda pozitif sınıf oranı eğitim dönemine göre doğrulama ve test dönemlerinde ciddi biçimde düşmektedir. Bu dağılım kayması F1 ve precision değerlerini sınırlamaktadır.
    
    \item \textbf{STFT öznitelikleri ayırt edici bilgi taşımaktadır:} Temiz normal/anomali ayrımında 51 boyutlu STFT öznitelikleriyle \%80 üzeri accuracy, precision, recall ve F1 değerleri elde edilmiştir.
    
    \item \textbf{Problem tanımı sonuçları belirleyicidir:} Deprem öncesi tahmin, olay çevresi anomali tespiti ve temiz normal/anomali ayrımı birbirinden farklı görevlerdir. Bu çalışmada en güçlü ve savunulabilir sonuç, temiz normal ve olay çevresi anomali pencerelerinin ayrılması görevinde elde edilmiştir.
\end{enumerate}

\subsection{Sınırlamalar}

Çalışmanın sınırlamaları aşağıda özetlenmiştir:

\begin{itemize}
    \item Sonuçlar ağırlıklı olarak tek istasyonun HHZ bileşeni üzerinden elde edilmiştir.
    \item Kronolojik testlerde pozitif sınıf oldukça seyrektir; bu durum F1 ve precision metriklerini doğrudan etkilemektedir.
    \item Temiz dengeli deney, olay çevresi ile olaylardan uzak normal dönemleri ayırır; bu nedenle deterministik deprem tahmini olarak yorumlanmamalıdır.
    \item Üç bileşenli veri kullanımı için HHZ, HHN ve HHE kanallarının zaman damgası seviyesinde hizalanması gerekmektedir.
\end{itemize}

\subsection{Genel Değerlendirme}

DeepSeis'in mevcut sürümü, sismik pencereleri STFT temsiline dönüştüren, farklı etiketleme protokollerini deneyen, kronolojik ve temiz dengeli değerlendirme sonuçlarını ayıran ve en başarılı modeli web arayüzünde çalıştırabilen uçtan uca bir prototip haline gelmiştir. En yüksek metrikler temiz normal/anomali ayrımı protokolünde elde edilmiştir. Bu nedenle tezde ana sonuç olarak bu protokol vurgulanmalı; kronolojik deprem öncesi tahmin sonuçları ise problemin zorluğunu ve gelecek çalışma ihtiyacını gösteren dürüst bir analiz olarak sunulmalıdır.

% -------------------- SONUÇ VE ÖNERİLER --------------------
\chapter{SONUÇ VE ÖNERİLER}

\section{Sonuçlar}

Bu tez çalışmasında DeepSeis adı verilen sismik anomali analizi prototipi geliştirilmiştir. Sistem; sürekli sismik kayıtların pencerelenmesi, STFT tabanlı zaman-frekans temsillerinin çıkarılması, farklı etiketleme stratejilerinin denenmesi, makine öğrenmesi ve derin öğrenme modellerinin karşılaştırılması ve seçilen modelin web arayüzünde çalıştırılması adımlarını kapsamaktadır.

Çalışmanın en önemli sonucu, problem tanımının model performansı üzerinde belirleyici olduğudur. Kronolojik bölümleme ile yapılan deprem öncesi kısa dönemli tahmin deneylerinde pozitif sınıf oranı test döneminde çok düşük kalmış ve F1 skorları sınırlı düzeyde gerçekleşmiştir. Buna karşın, temiz normal/anomali ayrımı protokolünde belirgin biçimde daha güçlü sonuçlar elde edilmiştir.

\subsection{Teknik Sonuçlar}

\begin{enumerate}
    \item \textbf{Temiz normal/anomali ayrımı:} Mw $\geq$ 4,0 olaylarının $\pm$60 dakika çevresi pozitif sınıf, olaylardan en az 48 saat uzak pencereler ise temiz normal sınıf olarak tanımlandığında HistGradientBoosting modeli 0,8124 accuracy, 0,8066 precision, 0,8219 recall ve 0,8142 F1 skoruna ulaşmıştır.
    
    \item \textbf{Skor tabanlı ayrım gücü:} Aynı modelde ROC-AUC 0,8979 ve PR-AUC 0,9090 olarak ölçülmüştür. Bu değerler, modelin normal ve olay çevresi anomali pencerelerini yalnızca tek bir eşikte değil, genel skor sıralaması bakımından da etkili biçimde ayırabildiğini göstermektedir.
    
    \item \textbf{Kronolojik testlerin zorluğu:} Zaman sırası korunan testlerde pozitif sınıf oranı yaklaşık \%1 seviyesine düşmüştür. Bu nedenle yüksek accuracy değerleri tek başına yeterli görülmemiş; precision, recall, F1, ROC-AUC, PR-AUC ve karışıklık matrisi birlikte değerlendirilmiştir.
    
    \item \textbf{STFT öznitelikleri:} STFT spektrogramlarından çıkarılan özet istatistiksel öznitelikler, temiz normal/anomali ayrımı görevinde güçlü bir temsil sağlamıştır. Bu bulgu, sismik sinyallerin frekans içeriğindeki değişimlerin anomali tespiti açısından önemli olduğunu desteklemektedir.
    
    \item \textbf{Uygulama prototipi:} En başarılı model, \texttt{joblib} formatında kaydedilmiş ve Flask tabanlı DeepSeis arayüzüne entegre edilmiştir. Arayüz, gelen test pencereleri için anomali olasılığı üretmekte ve karar eşiğine göre görsel uyarı sunmaktadır.
\end{enumerate}

\subsection{Bilimsel Yorum}

Bu çalışma deterministik deprem tahmini iddiası taşımamaktadır. Elde edilen güçlü sonuçlar, deprem olayına yakın sismik pencereler ile olaylardan uzak temiz normal pencereler arasında öğrenilebilir zaman-frekans farkları bulunduğunu göstermektedir. Bu nedenle sonuçlar, ``belirli koşullar altında sismik normal/anomali ayrımı'' olarak yorumlanmalıdır. Depremden saatler önce güvenilir uyarı üretmek ise daha uzun süreli, çok istasyonlu ve daha dengeli olay dağılımına sahip veri setleriyle ayrıca araştırılmalıdır.

\section{Sınırlılıklar ve Etik Uyarılar}

\begin{enumerate}
    \item \textbf{Tahmin iddiası:} DeepSeis çıktıları ``kesin deprem olacak'' şeklinde yorumlanmamalıdır. Sistem araştırma amaçlı bir anomali analiz prototipidir.
    
    \item \textbf{Tek istasyon ve kanal sınırı:} Mevcut güçlü sonuçlar HHZ bileşeni üzerinden elde edilmiştir. Çok bileşenli analiz için HHZ, HHN ve HHE kanallarının zaman damgası seviyesinde hizalanması gerekmektedir.
    
    \item \textbf{Temiz deney protokolü:} En iyi sonuçlar olay çevresi pencereleri ile olaylardan en az 48 saat uzak normal pencereler arasındaki kontrollü ayrımda elde edilmiştir. Bu protokol gerçek zamanlı geleceğe dönük tahmin senaryosundan farklıdır.
    
    \item \textbf{Veri dağılımı:} Deprem olayları veri zaman aralığı boyunca homojen dağılmamaktadır. Bu durum kronolojik testlerde train/test dağılım kaymasına neden olmaktadır.
    
    \item \textbf{Gürültü kaynakları:} Trafik, inşaat, hava koşulları ve yerel titreşimler gibi tektonik olmayan kaynaklar model skorlarını etkileyebilir.
\end{enumerate}

\section{Öneriler}

\subsection{Gelecek Araştırmalar}

\begin{enumerate}
    \item \textbf{Çoklu istasyon entegrasyonu:} Farklı istasyonların birlikte değerlendirilmesi, tek istasyon gürültüsünü azaltabilir ve mekansal örüntülerin öğrenilmesini sağlayabilir.
    
    \item \textbf{Üç bileşenli modelleme:} HHZ, HHN ve HHE kanallarının zaman hizalı biçimde birleştirilmesi, modelin polarizasyon ve yön bilgilerini kullanmasına olanak tanıyabilir.
    
    \item \textbf{Daha uzun veri aralığı:} Birkaç ay yerine yıllara yayılan veri kullanımı, daha fazla büyük deprem örneği sağlayarak kronolojik genelleme problemini daha sağlıklı değerlendirmeye yardımcı olacaktır.
    
    \item \textbf{Açıklanabilirlik analizi:} SHAP, permutation importance veya attention tabanlı yöntemlerle hangi frekans ve zaman bölgelerinin model kararında etkili olduğu incelenmelidir.
    
    \item \textbf{Gerçek zamanlı akış:} Mevcut arayüz test pencereleri üzerinde çalışmaktadır. Gelecek aşamada canlı MiniSEED akışı veya istasyon API'si üzerinden gerçek zamanlı veri beslemesi eklenebilir.
\end{enumerate}

\subsection{Uygulama Önerileri}

\begin{enumerate}
    \item Model çıktıları tek başına afet uyarısı olarak değil, uzman değerlendirmesine yardımcı araştırma sinyali olarak kullanılmalıdır.
    \item Arayüzde gösterilen anomali olasılığı, model protokolü ve karar eşiğiyle birlikte raporlanmalıdır.
    \item Yeni veri eklendikçe model performansı yeniden ölçülmeli ve eski metrikler güncel veri üzerinde doğrulanmalıdır.
\end{enumerate}

\section{Son Söz}

DeepSeis, bu çalışma sonunda yalnızca eğitim scriptlerinden oluşan bir deneme projesi olmaktan çıkmış; veri işleme, deney yönetimi, metrik raporlama, model kaydetme ve web arayüzünde model çalıştırma adımlarını içeren bütünlüklü bir prototipe dönüşmüştür. En başarılı deneyde elde edilen \%80 üzeri accuracy, precision, recall ve F1 değerleri, sismik zaman-frekans örüntülerinden temiz normal/anomali ayrımı yapılabileceğini göstermektedir. Bununla birlikte, deprem tahmini bilimsel olarak çok daha zor bir problem olduğundan, çalışmanın sonuçları dikkatli ve etik biçimde yorumlanmalıdır.

% -------------------- KAYNAKLAR --------------------
\begin{thebibliography}{99}
\addcontentsline{toc}{chapter}{KAYNAKLAR}

\bibitem{allen1978automatic}
Allen, R. V. (1978). Automatic earthquake recognition and timing from single traces. \textit{Bulletin of the Seismological Society of America}, 68(5), 1521--1532.

\bibitem{allen2009real}
Allen, R. M., Gasparini, P., Kamigaichi, O., \& Böse, M. (2009). The status of earthquake early warning around the world: An introductory overview. \textit{Seismological Research Letters}, 80(5), 682--693.

\bibitem{ambraseys2002seismicity}
Ambraseys, N. (2002). The seismic activity of the Marmara Sea region over the last 2000 years. \textit{Bulletin of the Seismological Society of America}, 92(1), 1--18.

\bibitem{bengio1994learning}
Bengio, Y., Simard, P., \& Frasconi, P. (1994). Learning long-term dependencies with gradient descent is difficult. \textit{IEEE Transactions on Neural Networks}, 5(2), 157--166.

\bibitem{erdik2003earthquake}
Erdik, M., Fahjan, Y., Ozel, O., Alcik, H., Mert, A., \& Gul, M. (2003). Istanbul earthquake rapid response and the early warning system. \textit{Bulletin of Earthquake Engineering}, 1(1), 157--163.

\bibitem{geller1997earthquakes}
Geller, R. J., Jackson, D. D., Kagan, Y. Y., \& Mulargia, F. (1997). Earthquakes cannot be predicted. \textit{Science}, 275(5306), 1616--1616.

\bibitem{given2014technical}
Given, D. D., Cochran, E. S., Heaton, T., Hauksson, E., Allen, R., Hellweg, P., \ldots \& Böse, M. (2014). Technical implementation plan for the ShakeAlert production system. \textit{U.S. Geological Survey Open-File Report}, 2014--1097.

\bibitem{hochreiter1997long}
Hochreiter, S., \& Schmidhuber, J. (1997). Long short-term memory. \textit{Neural Computation}, 9(8), 1735--1780.

\bibitem{kanamori2005real}
Kanamori, H. (2005). Real-time seismology and earthquake damage mitigation. \textit{Annual Review of Earth and Planetary Sciences}, 33, 195--214.

\bibitem{kong2019machine}
Kong, Q., Trugman, D. T., Ross, Z. E., Bianco, M. J., Meade, B. J., \& Gerstoft, P. (2019). Machine learning in seismology: Turning data into insights. \textit{Seismological Research Letters}, 90(1), 3--14.

\bibitem{li2019enhancing}
Li, S., Jin, X., Xuan, Y., Zhou, X., Chen, W., Wang, Y. X., \& Yan, X. (2019). Enhancing the locality and breaking the memory bottleneck of transformer on time series forecasting. \textit{Advances in Neural Information Processing Systems}, 32.

\bibitem{lim2021temporal}
Lim, B., Arık, S. Ö., Loeff, N., \& Pfister, T. (2021). Temporal fusion transformers for interpretable multi-horizon time series forecasting. \textit{International Journal of Forecasting}, 37(4), 1748--1764.

\bibitem{oppenheim1999discrete}
Oppenheim, A. V., Schafer, R. W., \& Buck, J. R. (1999). \textit{Discrete-time signal processing}. Prentice Hall.

\bibitem{perol2018convolutional}
Perol, T., Gharbi, M., \& Denolle, M. (2018). Convolutional neural network for earthquake detection and location. \textit{Science Advances}, 4(2), e1700578.

\bibitem{ross2018generalized}
Ross, Z. E., Meier, M. A., Hauksson, E., \& Heaton, T. H. (2018). Generalized seismic phase detection with deep learning. \textit{Bulletin of the Seismological Society of America}, 108(5A), 2894--2901.

\bibitem{sakurada2014anomaly}
Sakurada, M., \& Yairi, T. (2014). Anomaly detection using autoencoders with nonlinear dimensionality reduction. \textit{Proceedings of the MLSDA 2014 2nd Workshop on Machine Learning for Sensory Data Analysis}, 4--11.

\bibitem{scholz1990mechanics}
Scholz, C. H. (1990). \textit{The mechanics of earthquakes and faulting}. Cambridge University Press.

\bibitem{stein2009introduction}
Stein, S., \& Wysession, M. (2009). \textit{An introduction to seismology, earthquakes, and earth structure}. John Wiley \& Sons.

\bibitem{vaswani2017attention}
Vaswani, A., Shazeer, N., Parmar, N., Uszkoreit, J., Jones, L., Gomez, A. N., \ldots \& Polosukhin, I. (2017). Attention is all you need. \textit{Advances in Neural Information Processing Systems}, 30.

\bibitem{wu2021autoformer}
Wu, H., Xu, J., Wang, J., \& Long, M. (2021). Autoformer: Decomposition transformers with auto-correlation for long-term series forecasting. \textit{Advances in Neural Information Processing Systems}, 34.

\bibitem{zhou2021informer}
Zhou, H., Zhang, S., Peng, J., Zhang, S., Li, J., Xiong, H., \& Zhang, W. (2021). Informer: Beyond efficient transformer for long sequence time-series forecasting. \textit{Proceedings of the AAAI Conference on Artificial Intelligence}, 35(12), 11106--11115.

\end{thebibliography}

% -------------------- ÖZGEÇMİŞ --------------------
\chapter*{ÖZGEÇMİŞ}
\addcontentsline{toc}{chapter}{ÖZGEÇMİŞ}

\section*{Kişisel Bilgiler}

\begin{adjustbox}{center,max width=\textwidth}
\begin{tabular}{ll}
\textbf{Adı Soyadı:} & Can Ahmedi Yaşar PARLAK \\
\textbf{Doğum Yeri ve Tarihi:} & \\
\textbf{E-posta:} & \\
\end{tabular}
\end{adjustbox}

\section*{Eğitim Bilgileri}

\begin{adjustbox}{center,max width=\textwidth}
\begin{tabular}{p{3cm}p{9.5cm}}
\textbf{2020 - 2026} & Recep Tayyip Erdoğan Üniversitesi, Mühendislik ve Mimarlık Fakültesi, Bilgisayar Mühendisliği Bölümü (Lisans) \\
\textbf{20XX - 20XX} & Lise Eğitimi \\
\end{tabular}
\end{adjustbox}

\section*{Projeler ve Araştırmalar}

\begin{itemize}
    \item \textbf{DeepSeis:} Makine Öğrenmesi Tabanlı Sismik Anomali Tespiti ve Olay Çevresi Davranış Analizi (Bitirme Tezi, 2025-2026)
\end{itemize}

\section*{Teknik Beceriler}

\begin{itemize}
    \item \textbf{Programlama Dilleri:} Python, C++, JavaScript
    \item \textbf{Derin Öğrenme:} PyTorch, TensorFlow, Keras
    \item \textbf{Veri Bilimi:} NumPy, Pandas, Scikit-learn, Matplotlib
    \item \textbf{Sismik Veri İşleme:} ObsPy
    \item \textbf{Versiyon Kontrolü:} Git, GitHub
\end{itemize}

\section*{Yabancı Dil}

\begin{adjustbox}{center,max width=\textwidth}
\begin{tabular}{ll}
\textbf{İngilizce:} & Orta Düzey \\
\end{tabular}
\end{adjustbox}

\section*{İlgi Alanları}

\begin{itemize}
    \item Yapay Zeka ve Derin Öğrenme
    \item Zaman Serisi Analizi
    \item Sismoloji ve Deprem Bilimi
    \item Sinyal İşleme
\end{itemize}


\end{document}
